\Subhead{3.5.1}{Умножение частоты и амплитудное детектирование узкополосного шума}

Нелинейный элемент $y=a_1x+a_2x^2+\cdots$ смешивает частоты; после него необходимо указать фильтр, выбирающий нужную полосу. Для $x=R\cos(\omega_0t+\phi)$ квадрат содержит постоянную/медленную составляющую $R^2/2$ и гармонику $R^2\cos(2\omega_0t+2\phi)/2$. Поэтому частотный удвоитель меняет и несущую, и статистику огибающей; простой перенос спектра без изменения формы неверен.

Пусть $z(t)$ — центрированный круговой гауссовский процесс с корреляцией $K(\tau)=\langle z(t+\tau)z^*(t)\rangle$. Огибающая выделенной $n$-й гармоники пропорциональна $z^n$. По спариванию гауссовских моментов
\begin{equation}
 \langle z^n(t+\tau)[z^*(t)]^n\rangle=n!K^n(\tau).
\end{equation}
Для нормированной выходной огибающей $y=z^n$ её спектр равен $n$-кратной свёртке спектра $z$ с множителем $n!/(2\pi)^{n-1}$; коэффициент реального нелинейного преобразователя учитывается отдельно. Для экспоненциальной $K$ время корреляции уменьшается в $n$ раз, для гауссовой — в $\sqrt n$ раз: универсального правила «полоса умножается ровно на $n$» нет.

Квадратичный детектор после низкочастотного фильтра выделяет $R^2/2$. При $\langle X^2\rangle=\langle Y^2\rangle=\sigma^2$ имеем $\langle R^2\rangle=2\sigma^2$ и
\begin{equation}
 \operatorname{Cov}(|z(t+\tau)|^2,|z(t)|^2)=|K(\tau)|^2.
\end{equation}
Идеальный амплитудный детектор выделяет $R$, а не $R^2$: $\langle R\rangle=\sigma\sqrt{\pi/2}$, $\operatorname{Var}R=(2-\pi/2)\sigma^2$. Эти два закона детектирования имеют разную статистику. Конечная полоса фильтра усредняет огибающую и уменьшает дисперсию, если время усреднения превышает время корреляции.
